
\textbf{Théorème 4.6.}
\uline{Soient} \(X\) \uline{un schéma quasi-compact ayant un faisceau ample}, \(d\) \uline{un entier} \(>0\) \uline{et} \(i:Y\to X\) \uline{une immersion régulière de codimension} \(d\). \uline{Alors}, \uline{pour tout} \(k\), \uline{on a} :
\[
 i_*\bigl((\Fil^kK^{\bullet}(Y))\otimes_{\mathbb Z}\mathbb Q\bigr)
 \subset
 (\Fil^{k+d}K^{\bullet}(X))\otimes_{\mathbb Z}\mathbb Q.
\]

Remarquons tout d'abord que \(i_*(K^{\bullet}(Y))\subset \Fil^1K^{\bullet}(X)\). En effet, \(X-Y\) est dense dans \(X\) : s'il n'en était pas ainsi, \(Y\) contiendrait un ouvert affine, et sur cet ouvert, l'idéal \(\underline{J}\) définissant \(Y\) serait un nilidéal, donc un idéal localement nilpotent puisqu'il est de type fini. Or ceci est impossible, puisque les \(\underline{J}^n/\underline{J}^{n+1}\) sont des \(\underline{O}_Y\)-Modules localement libres de rang \(\ne0\). L'augmentation dans \(K^{\bullet}(X)\) d'un élément de \(i_*(K^{\bullet}(Y))\) est une fonction localement constante sur \(X\); étant nulle sur \(X-Y\), elle est nulle sur \(X\) entier.

Pour tout groupe abélien \(A\), nous poserons \(A_{\mathbb Q}=A\otimes_{\mathbb Z}\mathbb Q\). Soit \(x\in\Fil^kK^{\bullet}(Y)_{\mathbb Q}\) ; on peut écrire \(x\) sous la forme :
\[
  x=\sum a\,\gamma^{i_1}(x_1)\cdots\gamma^{i_n}(x_n),
\]
avec \(\epsilon(x_i)=0\) pour tout \(i\), \(a\in\mathbb Q\), et \(i_1+\cdots+i_n\ge k\). Les \(x_i\) peuvent s'exprimer comme combinaison à coefficients entiers des classes \(E_\alpha\) d'un nombre fini de \(\underline{O}_Y\)-Modules localement libres de type fini \(\underline{E}_\alpha\). Soit \(R\) la sous-\(\mathbb Q\)-\(\lambda\)-algèbre de \(K^{\bullet}(Y)_{\mathbb Q}\) engendrée par les \(E_\alpha\) et par \(N\). On a alors :
\[
  x\in\Fil^kR.
\]
Par ailleurs, la filtration de \(\lambda\)-anneau de \(R\) est discrète (VI 6.2); soit \(m\) tel que \(\Fil^mR=0\). Enfin, on peut, quitte à décomposer \(X\) en une somme finie d'ouverts disjoints et à faire la démonstration dans chaque ouvert (le \(K^{\bullet}\) commutant aux sommes directes finies), supposer les \(\underline{E}_\alpha\) de rang constant \(r_\alpha\). L'anneau d'augmentation de \(R\) est alors \(\mathbb Q\).

